\section{Usos}
% Cifras: tesis, cap. 5 (fuente de verdad; corrida 2026-10-04). Ver NOTAS_CORRECCION_ARTICULO.md
% ── Frame: Tres capas ────────────────────────────────────────────────────────
\begin{frame}{Apimetro como motor: analizar, crear, hacer}
  \centering
  \resizebox{!}{0.66\textheight}{% Diagrama: Macro-Arquitectura del Stack de Investigación (3 capas apiladas)
% Usado en: 4-Capas-Codigo/4-1-Apimetro/4-1-0-Filosofia.tex  (fig:macro-arquitectura-stack)
% Fix 2026-08-28: font=\small en estilo capa para que \hfill funcione; más espaciado vertical
% Copia para la presentación SOTM (26 oct 2026) — no modifica la tesis:
%   · Capa 1: ETL en Python + API en Go; Capa 3: QGIS + Tableau
%   · Cifras del grafo en la Capa 2 (tesis 5.2)
%   · Caja punteada "Capa N — tu proyecto" (puente al bloque 6)
%   · Capas de 12 cm (en la tesis 11 cm) para que quepa el título de la Capa 1
%   · Sin \gtfs{} ni \api{} (macros de la tesis, no definidos en el tema)
\tikzsetnextfilename{macro-arquitectura-stack}
\begin{tikzpicture}[
    capa/.style = {
        rectangle, draw=unamAzul!60,
        text width=12cm, minimum height=2.0cm,
        align=left, inner sep=14pt,
        rounded corners=4pt, font=\small
    },
    propia/.style = {
        rectangle, draw=SotmTerracota, dashed, line width=1pt,
        fill=SotmTerracota!6, text width=4.2cm,
        align=center, inner sep=12pt,
        rounded corners=4pt, font=\small
    },
    flecha/.style = {Stealth-Stealth, thick, color=unamAzul!55},
    etiq/.style   = {font=\scriptsize\itshape, color=gray!70,
                     align=center, text width=2.2cm}
]

%% — Capas (base más oscura, tope más clara) —
\node[capa, fill=unamAzul!28] (c1) {
    \textbf{Capa 1 — Apimetro: el motor de datos} \hfill \texttt{Python · Go · PostGIS}\\[5pt]
    ETL de datos abiertos + OSM · API REST · GeoJSON
};
\node[capa, fill=unamAzul!16, above=1.1cm of c1] (c2) {
    \textbf{Capa 2 — VFTModel: analizar} \hfill \texttt{Python · NetworkX}\\[5pt]
    Grafo de 11\,115 nodos y 45\,517 aristas · Indicadores topológicos
};
\node[capa, fill=unamAzul!7,  above=1.1cm of c2] (c3) {
    \textbf{Capa 3 — Transport-gis-zmvm-mjg: crear} \hfill \texttt{QGIS · Tableau}\\[5pt]
    Capas GeoPackage · Mapas · Tableros
};

%% — Flechas entre capas con etiquetas de intercambio —
\draw[flecha]
    ($(c1.north) + (0.9, 0)$) -- node[etiq, left=4pt] {GeoJSON}
    ($(c2.south) + (0.9, 0)$);

\draw[flecha]
    ($(c2.north) + (0.9, 0)$) -- node[etiq, left=4pt] {grafo /\\ resultados}
    ($(c3.south) + (0.9, 0)$);

%% — Flecha lateral "FLUJO DE DATOS" —
\draw[-Stealth, thick, color=unamAzul!50]
    ($(c1.south east) + (0.55, 0)$) -- ($(c3.north east) + (0.55, 0)$)
    node[midway, right=4pt,
         font=\footnotesize\bfseries\color{unamAzul!65},
         rotate=90, anchor=south] {FLUJO DE DATOS};

%% — Capa abierta: cualquier proyecto puede consumir la Capa 1 —
\node[propia, right=2.0cm of c2] (cn) {
    \textbf{\color{SotmTerracota}Capa N — tu proyecto}\\[5pt]
    Análisis, mapas, aplicaciones…\\ como tú quieras
};
\draw[-Stealth, thick, dashed, color=SotmTerracota]
    (c1.east) -| node[pos=0.75, right, font=\scriptsize\itshape,
                      color=SotmTerracota] {API REST} (cn.south);

\end{tikzpicture}
}\par
  \fuentefigura{Elaboración propia, 2026.}

  \vspace{0.15cm}
  % EDITAR: frase puente al bloque 6
  {\small\itshape\color{SotmVerde} VFTModel es solo un ejemplo: la Capa N puede ser la tuya.}
\end{frame}

% ── Frame: Cobertura (mapa QGIS) ─────────────────────────────────────────────
\begin{frame}{Cobertura: ¿quién tiene transporte a 800 m?}
  \begin{columns}[c]
    \begin{column}{0.66\textwidth}
      \centering
      \includegraphics[width=\linewidth,height=0.72\textheight,keepaspectratio]{Mapa_Cobertura_Mapa_Calor}\par
      \fuentefigura{Elaboración propia en QGIS con VFTModel \citep{cabrera2026VFTModel}.}
    \end{column}
    \begin{column}{0.32\textwidth}
      \centering
      {\small Milpa Alta\par}
      {\fontsize{28}{32}\selectfont\bfseries\color{SotmTerracota} 11.75\,\%\par}
      {\scriptsize de su superficie está a 800 m o menos de una estación\par}
      \vspace{0.3cm}
      {\small Alcaldías centrales: más del\par}
      {\fontsize{28}{32}\selectfont\bfseries\color{SotmVerde} 90\,\%\par}
      \vspace{0.3cm}
      % EDITAR
      {\footnotesize\itshape La desigualdad no es un accidente: es estructural.\par}
    \end{column}
  \end{columns}
\end{frame}

% ── Frame: Factor de desviación y fuerza capilar (tablero Tableau) ──────────
\begin{frame}{El costo de una red radial}
  \begin{columns}[c]
    \begin{column}{0.60\textwidth}
      \centering
      \includegraphics[width=\linewidth,height=0.72\textheight,keepaspectratio]{Dashboard_VFTModel_2026_Factor_desviacion}\par
      \fuentefigura{Elaboración propia en Tableau con VFTModel \citep{cabrera2026VFTModel}.}
    \end{column}
    \begin{column}{0.38\textwidth}
      \vspace*{-0.2cm}
      % EDITAR: lectura de las cifras
      \begin{exampleblock}{\footnotesize Factor de desviación (DI)}
        \scriptsize Media de \textbf{1.701}: cada viaje recorre cerca de
        \textbf{70\,\% más} que la línea recta (100 pares O-D).
      \end{exampleblock}
      \begin{block}{\footnotesize Caso extremo: DI = 2.99}
        \scriptsize Milpa Alta → Santa Catarina Yecahuizotl: 14.8 km en línea
        recta, \textbf{44.3 km} por la red.
      \end{block}
      \begin{alertblock}{\footnotesize Fuerza capilar: $FC_H = 130$}
        \scriptsize Hospital Gea González, Acoxpa (111) y La Rosa (100), sobre la
        Calzada de Tlalpan. Ningún gran CETRAM está entre los primeros.
      \end{alertblock}
    \end{column}
  \end{columns}
\end{frame}
